Este trabajo práctico nos permitió retomar lo aprendido en los trabajos
anteriores y aplicarlo a un sistema más completo. La ALU que construimos en el
TP1 volvió a formar parte del camino de datos, mientras que la comunicación
UART del TP2 nos permitió intercambiar programas y resultados con la PC. Esta
vez, el desafío no consistió en realizar una operación aislada, sino en integrar
esos bloques en un procesador capaz de ejecutar programas y permitirnos observar
su funcionamiento.

Al organizar el procesador en cinco etapas comprendimos que superponer
instrucciones mejora el aprovechamiento del circuito, pero también introduce
situaciones que tuvimos que resolver cuidadosamente. Una instrucción puede
necesitar un resultado que todavía no llegó al banco de registros y, al mismo
tiempo, un salto puede cambiar el camino cuando ya se buscaron instrucciones
posteriores. Trabajar con forwarding, bypass y esperas por \code{load-use} nos
ayudó a entender cuándo un dato está disponible y cómo hacer que cada
instrucción utilice la versión correcta.

También aprendimos a prestar atención a qué significa avanzar un ciclo del
procesador. Para poder ejecutar de manera continua o paso a paso y, a la vez,
observar el estado por UART, tuvimos que distinguir el reloj físico de los
ciclos efectivos de la CPU. El tratamiento de \code{halt} y los faults agregó
otra decisión importante: dejar terminar las instrucciones anteriores y
descartar las posteriores, de manera que el estado final conserve los efectos
que corresponden.

La carga y la preparación de programas nos mostraron la importancia de definir
claramente las responsabilidades de cada parte del sistema. El assembler
traduce el código, la UART transporta la imagen y el procesador la recibe y
prepara el estado inicial. Esto nos permitió ejecutar distintos programas sobre
una misma configuración de FPGA y volver a empezar una ejecución sin depender
de los datos que hubiera dejado la anterior. A su vez, Debug nos dio una forma
de relacionar lo que escribimos en el programa con lo que ocurre dentro del
pipeline.

Durante la verificación volvimos a apoyarnos en la metodología de los trabajos
anteriores: comprobar primero los bloques por separado y después observar cómo
funcionan al conectarlos. Las simulaciones nos ayudaron a seguir las
instrucciones ciclo a ciclo y a comprobar condiciones de borde; luego, la
implementación y las pruebas en la Basys 3 nos permitieron recorrer el sistema
con la placa y la comunicación reales.

En conjunto, el trabajo nos ayudó a comprender que diseñar un procesador
implica mucho más que implementar un conjunto de instrucciones. Tuvimos que
coordinar el tiempo y el estado de sus componentes, conectarlos con el exterior
y encontrar formas de verificar lo que ocurría en cada etapa. El resultado no
fue solamente un circuito capaz de ejecutar programas, sino también una
oportunidad para integrar los conceptos de la materia y llevarlos desde el
diseño hasta su observación en la FPGA.
